% =====================================================================
%  Chapter 30 — Assembling the System
% =====================================================================

\chapter{Assembling the System}
\label{ch:assembled}

This chapter describes the assembled system end to end: the monthly loop,
the order in which the layers act, and the accounting that runs alongside.
It is the chapter in which the reader can check that the architecture of
Parts IV and V is what the code actually does.

\section{The monthly loop}

At each month end \(t\), the system does the following, in this order.

\begin{enumerate}
\item \textbf{Receive dividends.} The cash dividends paid during month
\(t\) are credited to the accounts, in proportion to the positions held at
the end of month \(t-1$. No dividend is received on a position bought
during month \(t\).
\item \textbf{Form the target weights.} The strategic allocation is scaled
to the funds that exist; the base weights inside each sleeve are
inverse-volatility; the model's tilt is computed, clipped, scaled and
renormalised; and the risk governor then applies the volatility scaler
and the drawdown brake.
\item \textbf{Trade.} The target is compared with the current holdings.
Trades smaller than a no-trade band are skipped. Sells happen before
buys, so that the cash raised by a sale is available for a purchase in the
same month rather than the next.
\item \textbf{Price and value.} The positions are revalued at the next
month's prices, and each account's value is recorded.
\item \textbf{Settle the year.} If the month is a December, the taxable
account's tax for the year is computed and paid; the tax-exempt account
settles with no payment.
\end{enumerate}

The order matters, and each choice has a reason.

Dividends are credited before trading, because a dividend paid in month
\(t\) on a position held at the end of month \(t-1\) is income the
investor actually received and can reinvest.

Sells precede buys, so that the spread paid on a reallocation is charged
once per security rather than twice, and so that the account is not left
holding cash it did not need to hold.

The no-trade band exists for two reasons that Chapter 8 and Chapter 10
give separately: it saves spread costs, and it reduces the number of
moments at which the investor is invited to look at the account.

\section{Why the loop is monthly}

The data is monthly, the target horizons are one and three months, and the
governor acts on a drawdown that is measured in months. Monthly is the
natural frequency of the whole system.

But the reason to state the frequency is not the implementation; it is
Chapter 8. The evaluation frequency of the portfolio is a behavioural
parameter with a large price. A system that rebalances monthly but whose
owner looks at it monthly has made the myopic loss aversion of
Benartzi and Thaler fully operative on a strategy whose whole point is a
long horizon.

This is why the book's reporting cadence is deliberately slower than its
rebalancing cadence, and why the drawdown brake exists in the form it
does: the system is built so that the owner's attention is not required
at monthly frequency, even though the system's own arithmetic runs there.

\begin{takeaway}[The reporting cadence]
The system rebalances monthly and reports quarterly. The gap between the
two is a deliberate application of Chapter 8: the psychic cost of
evaluating a long-horizon portfolio more often than necessary is larger
than most of the effects this book measures.
\end{takeaway}

\section{The no-trade band and the cost model}

The no-trade band is half a per cent of the portfolio's value. A target
weight that differs from the current weight by less than that is not
traded. The effect on turnover is substantial, and the effect on the
tax bill is larger, because most of what the band suppresses is the
realisation of small gains.

The cost model charges half the quoted spread on each one-way trade. For
the most liquid funds the book charges two basis points of spread, for
the rest six. The costs are charged in the simulation, not estimated
afterwards, and Chapter 33 reports the result under doubled and tripled
spreads.

Expense ratios are not charged, for the reason set out in Chapter 10: they
are already inside the price series from which the returns are built.

\section{The two accounts}

The taxable and tax-exempt accounts are run in the same loop, with the same
targets and the same trades, and the only difference is that one pays tax
at each year end. They are also linked for the wash-sale rule: a
purchase in either account within the thirty-day window of a loss in the
taxable account disallows that loss.

The point of running both is that it isolates the tax exactly. A reader
who wants to know what the strategy is worth in a taxable account gets
the answer without having to model the tax themselves, and a reader who
wants to know what the strategy is worth in a retirement account gets
that answer too, from the same weights.

The governor's two traces over the record --- the exposure multiplier it applied and the drawdown that triggered the brake --- are shown in \figref{fig:governor}.

\section{What the assembly does not include}

The system's own record is a simulation. It is not a trading record, and
the book is careful never to describe it as one. It uses prices at month
end, charges a spread approximation rather than modelling the order
book, assumes the trades fill, assumes the investor can act at the month
end without market impact --- which Chapter 11 showed is safe at this
size --- and assumes no cash drag and no forced sales.

It also assumes the investor behaves: no overriding, no panic selling, no
additional purchases. Chapter 8's argument is that the system's design
is what makes that assumption reasonable, and Chapter 36 takes the
assumption seriously and asks what happens when it fails.

\begin{keytakeaways}
\item The monthly loop credits dividends, forms the target, trades with a
no-trade band and sells before buying, revalues, and settles the tax at
the year end.
\item The rebalancing cadence is monthly and the reporting cadence is
quarterly, which is a direct application of the myopic loss aversion
result.
\item The cost model charges half the spread per one-way trade and does
not charge expense ratios, because they are already in the returns.
\item The record is a simulation under stated assumptions, not a trading
record, and the book never describes it as one.
\end{keytakeaways}
